%% ============================================================
\chapter{Theoretical Analysis of the Fractional Model}
%% ============================================================

\begin{remarkbox}[title={Purpose of this chapter}]
This chapter establishes the main mathematical properties of the proposed fractional-order SITA model. The analysis verifies that the model is epidemiologically meaningful, that solutions remain nonnegative and bounded, and that the disease-free equilibrium can be studied through the basic reproduction number and fractional stability theory.
\end{remarkbox}

Throughout this chapter, the intervention functions are assumed to satisfy
\[
0\leq u_i(t)\leq 1,\qquad i=1,2,3,4,
\]
and all biological parameters are assumed to be nonnegative, with
\[
\Lambda>0,\quad \mu>0,\quad \beta_0>0,\quad \tau>0,\quad \delta>0,\quad \rho>0,\quad d>0,\quad 0<\eta<1.
\]
The fractional order satisfies
\[
0<q\leq 1.
\]

For mathematical analysis, the system is written in compact form as
\[
\CaputoD{t}{q}X(t)=F(t,X(t)),
\qquad
X(t)=\bigl(S(t),I(t),T(t),A(t)\bigr)^T.
\]

%% ------------------------------------------------------------
\section{Equivalent Integral Form}
%% ------------------------------------------------------------

The Caputo fractional system can be rewritten as an equivalent Volterra integral system by applying the Riemann--Liouville fractional integral. This representation is important because fractional differential equations are naturally nonlocal, meaning that the present state depends on the entire previous history of the system.

For the susceptible compartment, the equivalent integral form is
\begin{equation}
S(t)=S_0+\frac{1}{\Gamma(q)}
\int_0^t (t-s)^{q-1}
\left[\Lambda-\lambda(s)S(s)-\mu S(s)\right]ds.
\label{eq:int_S}
\end{equation}

Similarly, for the infected, treated, and AIDS-stage compartments, we obtain
\begin{align}
I(t)&=I_0+\frac{1}{\Gamma(q)}
\int_0^t (t-s)^{q-1}
\left[
\lambda(s)S(s)-\bigl(\tau(1+u_3(s))+\delta+\mu\bigr)I(s)
\right]ds,
\label{eq:int_I}\\[4pt]
T(t)&=T_0+\frac{1}{\Gamma(q)}
\int_0^t (t-s)^{q-1}
\left[
\tau(1+u_3(s))I(s)-\bigl(\rho(1-u_4(s))+\mu\bigr)T(s)
\right]ds,
\label{eq:int_T}\\[4pt]
A(t)&=A_0+\frac{1}{\Gamma(q)}
\int_0^t (t-s)^{q-1}
\left[
\delta I(s)+\rho(1-u_4(s))T(s)-(\mu+d)A(s)
\right]ds.
\label{eq:int_A}
\end{align}

\begin{remarkbox}[title={Interpretation}]
Equations~\eqref{eq:int_S}--\eqref{eq:int_A} show explicitly how the Gamma function enters the model. The kernel $(t-s)^{q-1}$ weights past states of the system, which is the mathematical source of memory in the fractional-order model.
\end{remarkbox}

%% ------------------------------------------------------------
\section{Existence and Uniqueness of Solutions}
%% ------------------------------------------------------------

For a mathematical model to be meaningful, it must have a solution and that solution should be unique under the stated initial conditions. Let
\[
\Omega_M=\left\{X\in\mathbb{R}_+^4: 0\leq S+I+T+A\leq M,\; S+I+T+A>0\right\}
\]
be a bounded positive region for some $M>0$. In this region, the force of infection
\[
\lambda(t)=\frac{\beta_0(1-u_1(t))(1-u_2(t))(I+\eta T)}{N}
\]
is continuous whenever $N>0$.

\begin{resultbox}[title={Theorem 4.1 --- Existence and Uniqueness}]
Suppose that the initial data satisfy
\[
S_0\geq0,\qquad I_0\geq0,\qquad T_0\geq0,\qquad A_0\geq0,
\qquad N_0=S_0+I_0+T_0+A_0>0.
\]
If the intervention functions $u_i(t)$, $i=1,2,3,4$, are bounded and continuous on $[0,T]$, then the fractional-order SITA system admits a unique local solution on a suitable interval $[0,T^*]\subseteq[0,T]$.
\end{resultbox}

\noindent\textbf{Proof.}
Using the Caputo integral representation, the system can be written as
\[
X(t)=X(0)+\frac{1}{\Gamma(q)}
\int_0^t(t-s)^{q-1}F(s,X(s))ds.
\]
Define the operator
\[
(\mathcal{T}X)(t)=X(0)+\frac{1}{\Gamma(q)}
\int_0^t(t-s)^{q-1}F(s,X(s))ds.
\]
On any bounded subset of $\Omega_M$, the function $F(t,X)$ is continuous in $t$ and locally Lipschitz in $X$, because it is made from sums, products, and rational terms whose denominator $N$ is positive. Hence, there exists a constant $L>0$ such that
\[
\|F(t,X)-F(t,Y)\|\leq L\|X-Y\|.
\]
For $X,Y\in C([0,T^*],\mathbb{R}^4)$, we have
\[
\|\mathcal{T}X-\mathcal{T}Y\|
\leq
\frac{L}{\Gamma(q)}
\int_0^{T^*}(T^*-s)^{q-1}\|X-Y\|ds.
\]
Therefore,
\[
\|\mathcal{T}X-\mathcal{T}Y\|
\leq
\frac{L(T^*)^q}{\Gamma(q+1)}\|X-Y\|.
\]
Choosing $T^*>0$ sufficiently small so that
\[
\frac{L(T^*)^q}{\Gamma(q+1)}<1,
\]
the operator $\mathcal{T}$ becomes a contraction. By Banach's fixed point theorem, the model has a unique local solution. Since boundedness is established below, the solution can be extended while it remains in the biologically feasible region. \hfill$\square$

%% ------------------------------------------------------------
\section{Positivity of Solutions}
%% ------------------------------------------------------------

Population variables cannot be negative. Thus, positivity is an essential biological requirement.

\begin{resultbox}[title={Theorem 4.2 --- Positivity}]
If
\[
S_0\geq0,\qquad I_0\geq0,\qquad T_0\geq0,\qquad A_0\geq0,
\]
then the solution of the fractional-order SITA model satisfies
\[
S(t)\geq0,\qquad I(t)\geq0,\qquad T(t)\geq0,\qquad A(t)\geq0,
\qquad t\geq0.
\]
\end{resultbox}

\noindent\textbf{Proof.}
We verify that the vector field points inward on the boundary of the nonnegative orthant. At $S=0$,
\[
\CaputoD{t}{q}S(t)=\Lambda\geq0.
\]
At $I=0$,
\[
\CaputoD{t}{q}I(t)=\lambda(t)S(t)\geq0.
\]
At $T=0$,
\[
\CaputoD{t}{q}T(t)=\tau(1+u_3(t))I(t)\geq0.
\]
At $A=0$,
\[
\CaputoD{t}{q}A(t)=\delta I(t)+\rho(1-u_4(t))T(t)\geq0.
\]
Since each boundary component has a nonnegative fractional derivative in the corresponding compartment, the nonnegative orthant is positively invariant by the standard positivity principle for Caputo fractional systems. Therefore, solutions starting with nonnegative initial conditions remain nonnegative for all $t\geq0$. \hfill$\square$

%% ------------------------------------------------------------
\section{Boundedness and Feasible Region}
%% ------------------------------------------------------------

Boundedness ensures that the total population remains finite. Let
\[
N(t)=S(t)+I(t)+T(t)+A(t).
\]
Adding the four model equations gives
\begin{equation}
\CaputoD{t}{q}N(t)=\Lambda-\mu N(t)-dA(t).
\label{eq:caputo_N}
\end{equation}
Since $dA(t)\geq0$, it follows that
\begin{equation}
\CaputoD{t}{q}N(t)\leq \Lambda-\mu N(t).
\label{eq:N_ineq}
\end{equation}

Consider the comparison equation
\[
\CaputoD{t}{q}Y(t)=\Lambda-\mu Y(t),\qquad Y(0)=N(0).
\]
Its solution is
\begin{equation}
Y(t)=\frac{\Lambda}{\mu}
+\left(N(0)-\frac{\Lambda}{\mu}\right)E_q(-\mu t^q),
\label{eq:Y_ML}
\end{equation}
where $E_q(\cdot)$ is the Mittag--Leffler function.

\begin{resultbox}[title={Theorem 4.3 --- Boundedness}]
All solutions of the fractional-order SITA model with nonnegative initial conditions are bounded in the region
\[
\Omega=
\left\{
(S,I,T,A)\in\mathbb{R}_+^4:
0\leq N(t)\leq \frac{\Lambda}{\mu}
\right\}
\]
whenever $N(0)\leq \Lambda/\mu$. In general,
\[
\limsup_{t\to\infty}N(t)\leq\frac{\Lambda}{\mu}.
\]
\end{resultbox}

\noindent\textbf{Proof.}
By the fractional comparison principle, inequality~\eqref{eq:N_ineq} implies
\[
N(t)\leq Y(t),
\]
where $Y(t)$ is given by~\eqref{eq:Y_ML}. Since $E_q(-\mu t^q)\to0$ as $t\to\infty$ for $\mu>0$, we obtain
\[
\limsup_{t\to\infty}N(t)\leq\frac{\Lambda}{\mu}.
\]
Therefore, the total population remains bounded and the region $\Omega$ is biologically feasible. \hfill$\square$

\begin{remarkbox}[title={Biological meaning}]
The bound $\Lambda/\mu$ represents the maximum long-term population size supported by recruitment and natural mortality in the absence of disease-induced mortality. The additional AIDS-induced death term $dA(t)$ only reduces the total population further.
\end{remarkbox}

%% ------------------------------------------------------------
\section{Disease-Free Equilibrium}
%% ------------------------------------------------------------

The disease-free equilibrium represents the state where HIV is absent from the population. Therefore,
\[
I=T=A=0.
\]
Substituting these values into the susceptible equation gives
\[
0=\Lambda-\mu S.
\]
Hence,
\begin{equation}
E_0=\left(\frac{\Lambda}{\mu},0,0,0\right).
\label{eq:dfe}
\end{equation}

\begin{remarkbox}[title={Interpretation}]
At $E_0$, all individuals are susceptible and no infection is present. This equilibrium is the reference point for deriving the basic reproduction number.
\end{remarkbox}

%% ------------------------------------------------------------
\section{Basic Reproduction Number}
%% ------------------------------------------------------------

For the derivation of $\Rzero$, the intervention functions are treated as constant levels:
\[
u_i(t)=u_i,\qquad i=1,2,3,4.
\]
Define
\[
\beff=\beta_0(1-u_1)(1-u_2),\qquad
\taueff=\tau(1+u_3),\qquad
\rhoeff=\rho(1-u_4).
\]
At the disease-free equilibrium,
\[
S^*=\frac{\Lambda}{\mu},
\qquad
N^*=\frac{\Lambda}{\mu}.
\]
Thus, near $E_0$,
\[
\lambda(t)S
=
\beff(I+\eta T).
\]

The infected-related compartments are $I$ and $T$. The new infection terms are
\[
\mathcal{F}_1=\beff(I+\eta T),
\qquad
\mathcal{F}_2=0.
\]
The transition terms are
\[
\mathcal{V}_1=(\taueff+\delta+\mu)I,
\]
and
\[
\mathcal{V}_2=(\rhoeff+\mu)T-\taueff I.
\]
Therefore,
\[
F=
\begin{pmatrix}
\beff & \eta\beff\\
0&0
\end{pmatrix},
\qquad
V=
\begin{pmatrix}
\taueff+\delta+\mu & 0\\
-\taueff & \rhoeff+\mu
\end{pmatrix}.
\]

The inverse of $V$ is
\[
V^{-1}=
\frac{1}{(\taueff+\delta+\mu)(\rhoeff+\mu)}
\begin{pmatrix}
\rhoeff+\mu & 0\\
\taueff & \taueff+\delta+\mu
\end{pmatrix}.
\]
The next-generation matrix is
\[
K=FV^{-1}.
\]
Multiplying gives
\[
K=
\frac{\beff}{(\taueff+\delta+\mu)(\rhoeff+\mu)}
\begin{pmatrix}
\rhoeff+\mu+\eta\taueff &
\eta(\taueff+\delta+\mu)\\
0&0
\end{pmatrix}.
\]
The basic reproduction number is the spectral radius of $K$. Hence,
\begin{mathbox}
\begin{equation}
\Rzero=
\frac{\beff}{\taueff+\delta+\mu}
\left(
1+\frac{\eta\taueff}{\rhoeff+\mu}
\right).
\label{eq:R0_general_eff}
\end{equation}
\end{mathbox}
Substituting the intervention-modified rates gives
\begin{mathbox}
\begin{equation}
\Rzero=
\frac{\beta_0(1-u_1)(1-u_2)}
{\tau(1+u_3)+\delta+\mu}
\left(
1+\frac{\eta\tau(1+u_3)}
{\rho(1-u_4)+\mu}
\right).
\label{eq:R0_intervention_final}
\end{equation}
\end{mathbox}

\begin{remarkbox}[title={Important interpretation of $u_4$}]
In the present model, $u_4$ reduces progression from treatment to AIDS through $\rho(1-u_4)$. Since the treated class still contributes to transmission through $\eta T$, the direct effect of $u_4$ on $\Rzero$ can be subtle. If the study aims to make adherence reduce infectiousness directly, a refined force of infection may use $\eta(1-u_4)T$. In the current structure, $u_4$ is mainly interpreted as reducing AIDS progression and improving treatment outcomes, while $u_1$ and $u_2$ directly reduce new infections.
\end{remarkbox}

%% ------------------------------------------------------------
\section{Local Stability of the Disease-Free Equilibrium}
%% ------------------------------------------------------------

The Jacobian matrix of the full system at $E_0$ is
\[
J(E_0)=
\begin{pmatrix}
-\mu & -\beff & -\eta\beff & 0\\
0 & \beff-(\taueff+\delta+\mu) & \eta\beff & 0\\
0 & \taueff & -(\rhoeff+\mu) & 0\\
0 & \delta & \rhoeff & -(\mu+d)
\end{pmatrix}.
\]
Two eigenvalues are immediately obtained:
\[
\lambda_1=-\mu,
\qquad
\lambda_2=-(\mu+d),
\]
which are both negative.

The remaining eigenvalues are obtained from the infected subsystem
\[
J_{IT}=
\begin{pmatrix}
\beff-(\taueff+\delta+\mu) & \eta\beff\\
\taueff & -(\rhoeff+\mu)
\end{pmatrix}.
\]
Let
\[
a=\taueff+\delta+\mu,
\qquad
b=\rhoeff+\mu.
\]
Then
\[
J_{IT}=
\begin{pmatrix}
\beff-a & \eta\beff\\
\taueff & -b
\end{pmatrix}.
\]
The trace is
\[
\operatorname{tr}(J_{IT})=\beff-a-b,
\]
and the determinant is
\[
\det(J_{IT})
=
(\beff-a)(-b)-\eta\beff\taueff.
\]
After simplification,
\begin{equation}
\det(J_{IT})
=
ab-\beff(b+\eta\taueff).
\label{eq:det_simplified}
\end{equation}
Using~\eqref{eq:R0_general_eff}, we obtain
\begin{equation}
\det(J_{IT})=ab(1-\Rzero).
\label{eq:det_R0_relation}
\end{equation}

\begin{resultbox}[title={Theorem 4.4 --- Fractional Local Stability of $E_0$}]
The disease-free equilibrium
\[
E_0=\left(\frac{\Lambda}{\mu},0,0,0\right)
\]
is locally asymptotically stable in the fractional-order sense if
\[
\Rzero<1
\]
and the eigenvalues $\lambda_i$ of the linearized system satisfy
\[
|\arg(\lambda_i)|>\frac{q\pi}{2}.
\]
It is unstable if $\Rzero>1$.
\end{resultbox}

\noindent\textbf{Proof.}
If $\Rzero<1$, then from~\eqref{eq:det_R0_relation},
\[
\det(J_{IT})=ab(1-\Rzero)>0.
\]
Also, $\Rzero<1$ implies
\[
\beff < \frac{ab}{b+\eta\taueff}<a,
\]
because $b+\eta\taueff>b$. Therefore,
\[
\operatorname{tr}(J_{IT})=\beff-a-b<0.
\]
Thus, the infected subsystem has eigenvalues with negative real parts. Since the remaining eigenvalues $-\mu$ and $-(\mu+d)$ are also negative, the linearized ordinary system is locally stable.

For a Caputo fractional system of order $0<q\leq1$, local asymptotic stability requires
\[
|\arg(\lambda_i)|>\frac{q\pi}{2}.
\]
Eigenvalues with negative real parts satisfy this condition for $0<q\leq1$. Hence $E_0$ is locally asymptotically stable in the fractional-order sense when $\Rzero<1$.

If $\Rzero>1$, then
\[
\det(J_{IT})=ab(1-\Rzero)<0,
\]
so the infected subsystem has eigenvalues with opposite signs in the real case or at least one eigenvalue with positive real part. Therefore, the disease-free equilibrium is unstable. \hfill$\square$

%% ------------------------------------------------------------
\section{Endemic Equilibrium}
%% ------------------------------------------------------------

An endemic equilibrium represents a persistent infection state:
\[
E^*=(S^*,I^*,T^*,A^*),
\qquad
I^*>0,\quad T^*>0,\quad A^*>0.
\]
At endemic equilibrium,
\[
\CaputoD{t}{q}S=
\CaputoD{t}{q}I=
\CaputoD{t}{q}T=
\CaputoD{t}{q}A=0.
\]
Therefore, the equilibrium equations are
\begin{align}
0&=\Lambda-\lambda^*S^*-\mu S^*,\\
0&=\lambda^*S^*-(\taueff+\delta+\mu)I^*,\\
0&=\taueff I^*-(\rhoeff+\mu)T^*,\\
0&=\delta I^*+\rhoeff T^*-(\mu+d)A^*.
\end{align}
From the third and fourth equations,
\[
T^*=\frac{\taueff}{\rhoeff+\mu}I^*,
\]
and
\[
A^*=
\frac{\delta I^*+\rhoeff T^*}{\mu+d}.
\]
The explicit expression for $I^*$ may be obtained by substituting these relations into the force of infection and susceptible equilibrium equation. Because the resulting expression is lengthy, this thesis uses numerical simulation to examine endemic behaviour when $\Rzero>1$.

\begin{remarkbox}[title={Interpretation}]
The disease-free equilibrium is the main analytical focus because it determines whether HIV can invade a fully susceptible population. The endemic equilibrium is mainly explored numerically, especially under scenarios where interventions are weak and $\Rzero>1$.
\end{remarkbox}

%% ------------------------------------------------------------
\section{Sensitivity Analysis}
%% ------------------------------------------------------------

The normalized sensitivity index of $\Rzero$ with respect to a parameter $p$ is
\begin{equation}
\Upsilon_p^{\Rzero}
=
\frac{\partial \Rzero}{\partial p}\cdot\frac{p}{\Rzero}.
\label{eq:sensitivity_index}
\end{equation}
This index measures the relative change in $\Rzero$ caused by a relative change in $p$.

From~\eqref{eq:R0_intervention_final}, the following direct sensitivity signs are obtained:
\[
\Upsilon_{\beta_0}^{\Rzero}=1,
\]
which means that a percentage increase in $\beta_0$ produces the same percentage increase in $\Rzero$.

For the awareness and safer behaviour interventions,
\[
\Upsilon_{u_1}^{\Rzero}=-\frac{u_1}{1-u_1},
\qquad
\Upsilon_{u_2}^{\Rzero}=-\frac{u_2}{1-u_2}.
\]
Thus, increasing $u_1$ or $u_2$ decreases $\Rzero$, confirming that awareness and safer sexual behaviour directly reduce transmission.

For the AIDS progression parameter $\delta$,
\[
\Upsilon_{\delta}^{\Rzero}
=
-\frac{\delta}{\taueff+\delta+\mu}.
\]
Hence increasing progression out of the untreated infected class reduces the average duration of untreated infection in this simplified transmission structure.

\begin{remarkbox}[title={Sensitivity interpretation}]
Positive sensitivity means that increasing the parameter increases $\Rzero$. Negative sensitivity means that increasing the parameter decreases $\Rzero$. The dashboard uses this idea to rank parameters and interventions according to their effect on HIV transmission potential.
\end{remarkbox}

%% ------------------------------------------------------------
\section{Chapter Summary}
%% ------------------------------------------------------------

This chapter established the main theoretical properties of the fractional-order SITA HIV model. The equivalent integral form showed how memory enters through the Gamma-weighted kernel. Existence and uniqueness were justified using a fixed-point argument, while positivity and boundedness confirmed the biological validity of the model. The disease-free equilibrium was obtained, the basic reproduction number was derived using the next-generation matrix method, and local stability was analysed using the fractional-order eigenvalue condition. Sensitivity analysis was introduced as a tool for identifying the most influential biological and intervention parameters.

%% ============================================================
\chapter{Research Methodology, Numerical Method and Dashboard Development}
%% ============================================================

\begin{remarkbox}[title={Purpose of this chapter}]
This chapter explains how the theoretical model is converted into numerical simulations and an interactive dashboard. It corrects and clarifies the fractional Adams--Bashforth--Moulton method, describes the computational workflow, and defines the simulation scenarios used to generate thesis results.
\end{remarkbox}

%% ------------------------------------------------------------
\section{Research Design}
%% ------------------------------------------------------------

This thesis uses a mathematical and computational research design. The mathematical component formulates and analyses the fractional-order SITA HIV model, while the computational component implements the model numerically and visualizes the results through an interactive dashboard.

The study follows three connected stages:
\begin{enumerate}
  \item \textbf{Model construction}: define compartments, intervention functions, fractional derivatives, and model assumptions.
  \item \textbf{Theoretical analysis}: establish positivity, boundedness, disease-free equilibrium, reproduction number, and fractional stability.
  \item \textbf{Computational simulation}: implement the fractional numerical method, run intervention scenarios, analyse memory effects, and present results through dashboard visualisations.
\end{enumerate}

\begin{remarkbox}[title={Nature of the results}]
The simulations in this thesis are scenario-based mathematical experiments. They are not direct clinical forecasts and are not fitted to patient-level data. Their purpose is to compare qualitative behaviour under different fractional orders and intervention levels.
\end{remarkbox}

%% ------------------------------------------------------------
\section{Model Development Procedure}
%% ------------------------------------------------------------

The development of the fractional-order HIV model follows the sequence below:
\begin{enumerate}
  \item Define the total population as
  \[
  N(t)=S(t)+I(t)+T(t)+A(t).
  \]
  \item Construct the SITA compartmental structure.
  \item Introduce four intervention functions:
  \[
  u_1,u_2,u_3,u_4\in[0,1].
  \]
  \item Define the effective rates:
  \[
  \beff=\beta_0(1-u_1)(1-u_2),\qquad
  \taueff=\tau(1+u_3),\qquad
  \rhoeff=\rho(1-u_4).
  \]
  \item Replace ordinary derivatives by the Caputo fractional derivative $\CaputoD{t}{q}$, where $0<q\leq1$.
  \item Derive the reproduction number $\Rzero$.
  \item Implement the fractional predictor--corrector method in Python.
  \item Use dashboard graphs and tables to compare baseline, intervention, and memory-effect scenarios.
\end{enumerate}

%% ------------------------------------------------------------
\section{Fractional Numerical Scheme}
%% ------------------------------------------------------------

The fractional-order model is solved numerically using the fractional Adams--Bashforth--Moulton predictor--corrector method. This method is widely used for Caputo fractional differential equations because it approximates the Volterra integral representation of the system and naturally accounts for memory effects.

Consider the general Caputo fractional initial value problem
\begin{equation}
\CaputoD{t}{q}y(t)=f(t,y(t)),
\qquad
y(0)=y_0,
\qquad
0<q\leq1.
\label{eq:general_fivp}
\end{equation}
Its equivalent integral form is
\begin{equation}
y(t)=y_0+\frac{1}{\Gamma(q)}
\int_0^t(t-s)^{q-1}f(s,y(s))ds.
\label{eq:general_integral}
\end{equation}
Let
\[
t_n=nh,\qquad n=0,1,\ldots,M,
\]
where $h>0$ is the time step.

%% ------------------------------------------------------------
\subsection{Predictor Step}
%% ------------------------------------------------------------

The Adams--Bashforth predictor estimates $y_{n+1}$ using previously computed values:
\begin{equation}
y_{n+1}^{P}
=
y_0+\frac{h^q}{\Gamma(q+1)}
\sum_{j=0}^{n} b_{j,n+1} f(t_j,y_j),
\label{eq:correct_predictor}
\end{equation}
where
\begin{equation}
b_{j,n+1}=(n+1-j)^q-(n-j)^q,
\qquad
j=0,1,\ldots,n.
\label{eq:correct_predictor_weights}
\end{equation}

%% ------------------------------------------------------------
\subsection{Corrector Step}
%% ------------------------------------------------------------

The Adams--Moulton corrector improves the predicted value by evaluating the function at $t_{n+1}$:
\begin{equation}
y_{n+1}
=
y_0+
\frac{h^q}{\Gamma(q+2)}
\left[
f(t_{n+1},y_{n+1}^{P})
+
\sum_{j=0}^{n}a_{j,n+1}f(t_j,y_j)
\right].
\label{eq:correct_corrector}
\end{equation}

The corrector weights are
\begin{equation}
a_{j,n+1}=
\begin{cases}
n^{q+1}-(n-q)(n+1)^q, & j=0,\\[4pt]
(n-j+2)^{q+1}+(n-j)^{q+1}-2(n-j+1)^{q+1}, & 1\leq j\leq n,\\[4pt]
1, & j=n+1.
\end{cases}
\label{eq:correct_corrector_weights}
\end{equation}

\begin{resultbox}[title={Corrected ABM weight formula}]
The first corrector coefficient is
\[
a_{0,n+1}=n^{q+1}-(n-q)(n+1)^q.
\]
This is the standard form used in the Diethelm fractional Adams--Bashforth--Moulton predictor--corrector method. It replaces the earlier indexing form that may lead to incorrect weights.
\end{resultbox}

%% ------------------------------------------------------------
\section{Application of the Scheme to the SITA System}
%% ------------------------------------------------------------

For the SITA system, define
\[
Y_n=(S_n,I_n,T_n,A_n)^T.
\]
The right-hand side vector is
\[
F(t_n,Y_n)=
\begin{pmatrix}
\Lambda-\lambda_nS_n-\mu S_n\\
\lambda_nS_n-(\tau(1+u_3)+\delta+\mu)I_n\\
\tau(1+u_3)I_n-(\rho(1-u_4)+\mu)T_n\\
\delta I_n+\rho(1-u_4)T_n-(\mu+d)A_n
\end{pmatrix},
\]
where
\[
\lambda_n=
\frac{\beta_0(1-u_1)(1-u_2)(I_n+\eta T_n)}
{S_n+I_n+T_n+A_n}.
\]

The vector predictor is
\begin{equation}
Y_{n+1}^{P}
=
Y_0+\frac{h^q}{\Gamma(q+1)}
\sum_{j=0}^{n}b_{j,n+1}F(t_j,Y_j),
\label{eq:vector_predictor}
\end{equation}
and the vector corrector is
\begin{equation}
Y_{n+1}
=
Y_0+\frac{h^q}{\Gamma(q+2)}
\left[
F(t_{n+1},Y_{n+1}^{P})
+
\sum_{j=0}^{n}a_{j,n+1}F(t_j,Y_j)
\right].
\label{eq:vector_corrector}
\end{equation}

After each numerical step, nonnegative projection is applied:
\[
Y_{n+1}\leftarrow \max(Y_{n+1},0),
\]
component by component. This prevents small numerical round-off errors from producing negative population values.

%% ------------------------------------------------------------
\section{Numerical Algorithm}
%% ------------------------------------------------------------

\begin{resultbox}[title={Algorithm: Fractional ABM Simulation for the SITA Model}]
\begin{enumerate}
  \item Set the initial vector
  \[
  Y_0=(S_0,I_0,T_0,A_0)^T.
  \]
  \item Select model parameters and intervention levels.
  \item Choose the fractional order $q$, final time $T_f$, and step size $h$.
  \item Generate the time grid $t_n=nh$, $n=0,1,\ldots,M$.
  \item For each time step $n=0,1,\ldots,M-1$:
  \begin{enumerate}
    \item Compute predictor weights $b_{j,n+1}$.
    \item Evaluate the predictor $Y_{n+1}^{P}$.
    \item Compute corrector weights $a_{j,n+1}$.
    \item Evaluate the corrected value $Y_{n+1}$.
    \item Apply nonnegative projection if necessary.
  \end{enumerate}
  \item Compute $\Rzero$, peak infection, final states, and summary indicators.
  \item Send results to the dashboard for visualisation.
\end{enumerate}
\end{resultbox}

%% ------------------------------------------------------------
\section{Numerical Accuracy and Consistency Checks}
%% ------------------------------------------------------------

The numerical method is checked using the following consistency principles:
\begin{enumerate}
  \item \textbf{Non-negativity}: all simulated compartments must remain nonnegative.
  \item \textbf{Boundedness}: the total population $N(t)$ should remain within the feasible region predicted by the theoretical analysis.
  \item \textbf{Classical limit}: when $q=1$, the fractional model should behave like the corresponding ordinary differential equation model.
  \item \textbf{Step-size reliability}: reducing the step size $h$ should not change the qualitative conclusions of the simulation.
  \item \textbf{Threshold consistency}: scenarios with $\Rzero<1$ should generally show declining infection, while scenarios with $\Rzero>1$ may show persistence.
\end{enumerate}

\begin{remarkbox}[title={Part to fill after simulation}]
After running final dashboard simulations, insert a short table comparing results for at least two step sizes, for example $h=0.2$ and $h=0.1$. This will strengthen the numerical credibility of the thesis.
\end{remarkbox}

%% ------------------------------------------------------------
\section{Simulation Scenarios}
%% ------------------------------------------------------------

Eight simulation scenarios are used to examine the effects of social behaviour interventions and fractional memory. The scenarios are designed for comparative mathematical analysis rather than direct clinical prediction.

\begin{table}[H]
\centering
\caption{Simulation scenarios used for intervention and memory comparison.}
\label{tab:simulation_scenarios_improved}
\small
\begin{tabularx}{\textwidth}{c l c c c c c X}
\toprule
No. & Scenario & $u_1$ & $u_2$ & $u_3$ & $u_4$ & $q$ & Purpose\\
\midrule
1 & No intervention & 0.0 & 0.0 & 0.0 & 0.0 & 0.95 & Baseline fractional trajectory without behavioural control.\\
2 & Awareness only & 0.7 & 0.0 & 0.0 & 0.0 & 0.95 & Isolates the effect of awareness and education.\\
3 & Safer behaviour only & 0.0 & 0.7 & 0.0 & 0.0 & 0.95 & Isolates the effect of condom use and safer behaviour.\\
4 & Testing only & 0.0 & 0.0 & 0.7 & 0.0 & 0.95 & Tests faster movement from infection to treatment.\\
5 & Adherence only & 0.0 & 0.0 & 0.0 & 0.8 & 0.95 & Tests reduced progression from treatment to AIDS.\\
6 & Combined intervention & 0.7 & 0.7 & 0.6 & 0.8 & 0.95 & Represents integrated prevention, testing, treatment, and adherence support.\\
7 & Ordinary model & 0.5 & 0.5 & 0.5 & 0.5 & 1.00 & Removes fractional memory for comparison with the classical model.\\
8 & High memory effect & 0.4 & 0.4 & 0.4 & 0.4 & 0.75 & Tests stronger memory effects on the epidemic trajectory.\\
\bottomrule
\end{tabularx}
\end{table}

\begin{scenariobox}[title={Scenario interpretation}]
Scenarios 1--5 identify individual intervention effects. Scenario 6 tests a combined strategy. Scenarios 7 and 8 isolate the impact of the fractional order by comparing the ordinary model with a stronger-memory fractional model.
\end{scenariobox}

%% ------------------------------------------------------------
\section{Fractional-Order Comparison}
%% ------------------------------------------------------------

To study memory effects, the system is simulated using
\[
q=1.00,\qquad q=0.95,\qquad q=0.85,\qquad q=0.75.
\]
When $q=1$, the model reduces to the classical first-order system. When $q<1$, memory effects are introduced through the Caputo kernel. Smaller values of $q$ represent stronger memory and slower decay of past influence.

The dashboard compares the infected curve $I(t)$, the treated curve $T(t)$, the AIDS-stage curve $A(t)$, and the $I$--$T$ phase trajectory under different values of $q$.

%% ------------------------------------------------------------
\section{Dashboard Architecture}
%% ------------------------------------------------------------

The dashboard is implemented as a small scientific web system with five layers:

\begin{center}
\renewcommand{\arraystretch}{1.25}
\begin{tabularx}{\textwidth}{>{\bfseries\color{urblue}}p{4cm}X}
\toprule
Layer & Function\\
\midrule
Model layer & Stores the fractional SITA equations and intervention functions.\\
Solver layer & Runs the fractional ABM predictor--corrector numerical method.\\
Analysis layer & Computes $\Rzero$, final states, peak infection, and sensitivity indices.\\
API layer & Provides simulation, scenario, memory comparison, sensitivity, and export endpoints.\\
Interface layer & Displays sliders, cards, graphs, tables, and downloadable outputs.\\
\bottomrule
\end{tabularx}
\end{center}

%% ------------------------------------------------------------
\section{Dashboard Features}
%% ------------------------------------------------------------

The dashboard contains the following major features:
\begin{itemize}
  \item Fractional order slider for $q$.
  \item Parameter sliders for biological and intervention parameters.
  \item Real-time $\Rzero$ computation.
  \item Time-series graphs for $S(t)$, $I(t)$, $T(t)$, $A(t)$, and $N(t)$.
  \item Scenario comparison module.
  \item Memory-effect comparison for different values of $q$.
  \item Sensitivity analysis bar chart.
  \item Dashboard screenshot and result export functionality.
  \item CSV export for simulation outputs.
\end{itemize}

%% ------------------------------------------------------------
\section{Suggested Project Folder Structure}
%% ------------------------------------------------------------

\begin{lstlisting}[style=pythonstyle]
fractional_hiv_thesis_dashboard/
|
|-- app.py
|-- config.py
|-- requirements.txt
|-- README.md
|
|-- model/
|   |-- fractional_sita_model.py
|   |-- fractional_solver.py
|   |-- reproduction_number.py
|   |-- scenarios.py
|   |-- sensitivity.py
|   `-- validation.py
|
|-- routes/
|   |-- simulation_routes.py
|   |-- export_routes.py
|   `-- health_routes.py
|
|-- templates/
|   |-- base.html
|   |-- index.html
|   |-- dashboard.html
|   |-- documentation.html
|   `-- about.html
|
|-- static/
|   |-- css/
|   |-- js/
|   |-- assets/
|   `-- vendor/
|
|-- data/
|   |-- baseline_parameters.json
|   |-- scenario_presets.json
|   `-- sample_outputs.csv
|
`-- outputs/
    |-- csv/
    |-- figures/
    `-- reports/
\end{lstlisting}

%% ------------------------------------------------------------
\section{Chapter Summary}
%% ------------------------------------------------------------

This chapter described the computational methodology used to simulate the fractional-order SITA HIV model. The fractional Adams--Bashforth--Moulton predictor--corrector scheme was presented with corrected weights, and its application to the four-dimensional SITA system was described. The chapter also explained the numerical algorithm, consistency checks, simulation scenarios, fractional-order comparison, and dashboard architecture. The next chapter presents the numerical results and discusses their epidemiological and computational interpretation.
